\section{Discusión}
A partir de la tabulación sistemática y las tendencias estadísticas obtenidas del corpus, resulta de gran interés contrastar cualitativamente las implicaciones prácticas de los modelos en seguridad vial y control de tránsito.

\subsection{Implicaciones del Monitoreo de Conductores y Tránsito}
En la categoría de monitoreo de conductores, la precisión de la detección de fatiga y expresiones de somnolencia se ha visto drásticamente incrementada gracias al uso de modelos basados en marcas faciales y redes convolucionales espaciotemporales. Trabajos como los recopilados en \cite{kim2025unsupervised, essahraui2025realtime} demuestran que el análisis ocular y bucal combinado previene efectivamente colisiones reales viales en vehículos comerciales. Por otro lado, la clasificación de emociones del conductor en tiempo real mediante deep learning \cite{lin2025dangerous, le2025redlight} permite estimar estados de ira, estrés o cansancio crónico que correlacionan directamente con una conducción agresiva o errática.

En lo que respecta a la detección de infracciones y monitoreo de tráfico, los estudios destacan la transición de detectores estáticos a sistemas dinámicos inteligentes de control. El uso de la familia YOLO en intersecciones urbanas complejas ha permitido rastrear y clasificar vehículos simultáneamente, identificando infracciones críticas como el giro indebido, el cruce de semáforo en rojo o el exceso de velocidad instantáneo mediante el cálculo de distancias relativas entre cuadros. Estudios en CCTV \cite{pavan2026aipowered, le2025redlight} y bajo tomas aéreas \cite{abukhadrah2025droneassisted, chouhan2025deep} recalcan que las arquitecturas modernas superan la latencia de procesamiento requerida, logrando tasas de precisión robustas en entornos metropolitanos de alta congestión.

\subsection{Inteligencia Artificial Explicable (XAI) en Seguridad Activa}
La falta de interpretabilidad en las decisiones de los modelos de Deep Learning (el problema de la "caja negra") representa una barrera significativa para la certificación industrial de sistemas de asistencia al conductor (ADAS) \cite{le2025redlight, silva2025light}. Ante una alerta de fatiga o distracción, es indispensable auditar si el modelo basó su predicción en patrones válidos (como la tasa de cierre de párpados o el bostezo) o en sesgos irrelevantes del fondo o del sujeto \cite{oussouaddi2025dsryolo, abukhadrah2025droneassisted}. Estudios recientes proponen el uso de técnicas de XAI visual, como Grad-CAM (Gradient-weighted Class Activation Mapping) \cite{deshpande2025computervision, elhenidy2025gyyolo}, para generar mapas de calor sobre el rostro que resaltan las regiones críticas de decisión en tiempo de inferencia, garantizando la auditabilidad y fomentando la confianza del usuario final.

\subsection{El Desafío 'In-the-Wild' y la Transferibilidad de Dominios}
Una limitación crítica en el estado del arte es el sesgo de los datasets de entrenamiento, los cuales suelen ser capturados en laboratorios con condiciones controladas de luz y pose. Al desplegar estos modelos en la conducción naturalista real ("in-the-wild"), se observa un decremento sustancial en el desempeño debido a la vibración del vehículo, reflejos del parabrisas y variabilidad climática extrema \cite{lu2024validation, khattak2024role}. Para superar esta limitación, las investigaciones se orientan al desarrollo de técnicas de Adaptación de Dominio No Supervisada (UDA) \cite{barbu2025deep, yadav2025enhancing} y al entrenamiento auto-supervisado, permitiendo que el sistema adapte su representación de características de un entorno sintético o controlado hacia la dinámica del mundo real sin requerir un costoso etiquetado manual adicional \cite{chouhan2025deep, singh2025realtime}.
